% ==============================================================================
% SECAO 7 — ANALISE DE COMPROMISSO E CONCLUSOES
% ==============================================================================
\section{ANÁLISE DE COMPROMISSO (\textit{TRADE-OFFS}) E CONCLUSÕES}\label{sec:conclusoes}

Esta seção discute os compromissos de engenharia revelados pela Tabela~\ref{tab:comparativo}, justifica a escolha do controlador e apresenta as considerações finais.

\subsection{Síntese Comparativa e Escolha do Controlador}\label{sec:escolha_otima}
A análise conjunta dos resultados analíticos, frequenciais e temporais indica que, entre as estruturas avaliadas, apenas a PID satisfaz as seis especificações da Seção~\ref{sec:requisitos}, pois o PI é limitado pela invariância de Vieta e o PD não rejeita a perturbação permanente.

Em ambiente puramente linear irrestrito, o PID nominal ($\alpha = 0{,}15\text{ s}^{-1}$, $K_p=2{,}160$, $K_i=0{,}243$, $K_d=2{,}900$) cumpre todas as metas ($t_r=1{,}022\text{ s}$, $t_s=10{,}440\text{ s}$, $M_p=8{,}13\%$, $e_{ss,\text{dist}}=0{,}15\text{ eventos}$, $\mathrm{PM}=77{,}67^\circ$ e $\mu_{\max}=1397{,}0\text{ eventos/s}$), embora à custa de pico teórico impraticável na capacidade solicitada $\mu_c(0^+)$. Já sob as restrições físicas reais da plataforma (saturação de $\mu_c$ em $[600,\,1400]\text{ eventos/s}$ e não negatividade da fila), o PID ajustado ($\alpha = 0{,}40\text{ s}^{-1}$, $K_p=3{,}060$, $K_i=0{,}648$, $K_d=3{,}400$), implementado com filtro derivativo ($\tau_f = 0{,}05\text{ s}$), derivada na medição e \textit{anti-windup} por \textit{back-calculation} ($T_t = 0{,}5\text{ s}$), atende a todos os critérios: tempo de subida $t_r = 2{,}370\text{ s} \le 5\text{ s}$, tempo de acomodação $t_s = 10{,}744\text{ s} \le 12\text{ s}$, sobressinal $M_p = 6{,}98\% \le 10\%$, erro nulo em regime permanente ($e_{ss} = 0{,}00\text{ eventos} \le 10\text{ eventos}$), margem de fase de $73{,}28^\circ \ge 45^\circ$ e capacidade efetiva contida no intervalo operacional ($\mu(t) \in [996{,}15,\,1130{,}21] \subset [600,\,1400]\text{ eventos/s}$).

\subsection{Compromisso entre Estabilidade, Rapidez, Precisão e Consumo de Recursos}\label{sec:tradeoffs}
O confronto entre as arquiteturas evidencia os compromissos estruturais inerentes ao controle dessa planta integradora com atraso de primeira ordem. O controlador P distingue-se pela simplicidade algorítmica e pelo baixo consumo de capacidade ($\mu_{\max}=1110{,}5\text{ eventos/s}$ para $K=0{,}369$), mas não atende aos requisitos devido à invariância da parte real dos polos ($\sigma = -0{,}25\text{ s}^{-1}$), que impõe acomodação lenta ($13{,}752\text{ s}$), e à incapacidade de absorver desvios permanentes de carga ($271{,}00\text{ eventos}$ de erro residual após a perturbação). A adição exclusiva de ação integral no controlador PI elimina o erro permanente de perturbação ($e_{ss,\text{dist}} \to 0$) sem gerar picos derivativos, porém a invariância de Vieta ($p_1+p_2+p_3=-0{,}5\text{ s}^{-1}$) impõe um conflito insolúvel entre lentidão e oscilação ($M_p = 34{,}67\%$, $t_s = 19{,}578\text{ s}$ e $\mathrm{PM} = 42{,}78^\circ$), além de introduzir limite finito de estabilidade ($K_i < K_p/2$).

Por sua vez, o controlador PD com derivada na medição alcança amortecimento elevado e rapidez na partida ($t_s = 4{,}762\text{ s}$, $M_p = 1{,}18\%$) com consumo reduzido de infraestrutura ($\mu_{\max} = 1104{,}7\text{ eventos/s}$), mas retém erro estacionário elevado após mudanças permanentes na chegada ($e_{ss,\text{dist}} = 78{,}13\text{ eventos}$) porque a perturbação incide a montante do integrador da planta. O PID com derivada na medição, filtro passa-baixas e \textit{back-calculation} concilia essas demandas antagônicas: a ação derivativa antecipa a tendência da fila e desloca a soma dos polos para a esquerda (assegurando $\mathrm{PM} = 73{,}28^\circ$ e $M_p = 6{,}98\%$); a ação integral garante erro nulo mesmo diante da sobrecarga permanente de $+100\text{ eventos/s}$; a derivada aplicada exclusivamente à medição elimina o \textit{derivative kick}, evitando escalonamentos bruscos de infraestrutura; e o \textit{anti-windup} impede que a saturação transitória ($0{,}534\text{ s}$) degrade o desempenho (ÅSTRÖM; HÄGGLUND, 2006). Soma-se a isso o fato de o pico efetivo de processamento atingir apenas $1130{,}21\text{ eventos/s}$, preservando uma margem de segurança próxima de $270\text{ eventos/s}$ abaixo do limite máximo instalado.

\subsection{Considerações Finais}\label{sec:consideracoes_finais}
Retomando o desafio inicial de prover elasticidade automática e garantia de SLA a plataformas de \textit{streaming} IoT descrito na Seção~\ref{sec:intro}, os resultados deste estudo confirmam que a modelagem física e matemática aliada à sintonia clássica oferece uma alternativa fundamentada às políticas empíricas de provisionamento, produzindo malhas de controle com desempenho previsível. Ao integrar conservação de fluxo, análise no plano-$s$ e na frequência e simulação temporal não linear, verificou-se que especificações aparentemente atendidas em modelo linear ideal, como no caso do PID nominal com $\alpha = 0{,}15\text{ s}^{-1}$, podem ser degradadas na prática caso a saturação de recursos e o \textit{windup} integral não sejam tratados no projeto.

Conclui-se que o PID ajustado ($\alpha = 0{,}40\text{ s}^{-1}$, $K_p=3{,}060$, $K_i=0{,}648$, $K_d=3{,}400$), com filtro derivativo ($\tau_f = 0{,}05\text{ s}$), derivada na medição e \textit{anti-windup} por \textit{back-calculation} ($T_t = 0{,}5\text{ s}$), é a configuração recomendada para a plataforma: atende a todas as especificações, recupera-se da sobrecarga permanente em $5{,}950\text{ s}$ sem erro residual e mantém a capacidade efetiva cerca de $270\text{ eventos/s}$ abaixo do limite instalado. Esses resultados pressupõem, contudo, um modelo contínuo e determinístico, com constante de tempo $\tau$ fixa, medição instantânea e sem ruído, e capacidade ajustável de forma contínua. A extensão natural do estudo inclui a análise de robustez frente a variações de $\tau$, a inclusão de atraso de medição e de discretização do ciclo de escalonamento e a consideração da granularidade das instâncias computacionais.
